\section{Prime graphs and simultaneous unit coefficients}\label{gr:section}

The coefficient symbols will now be regarded as functions of finitely
many quadratic residue bits. We first identify the arithmetic data that
those bits record and explain how the trace calculation becomes a
finite difference. This gives polynomials whose derivatives contract
prime labels. We then adjoin auxiliary vertices so that one or two
nonzero coefficient polynomials take the value $1$ at every nonzero
address of a binary cube.

The two-polynomial case is needed for the odd converse: the odd symbol
tests $h$, while an even symbol tests a quadratic companion $kh$.
Separate nonzero polynomials need not have a common unit evaluation;
for example, $z$ and $1+z$ do not. The auxiliary vertices change the
polynomials by controlled contraction operators. The arithmetic input
is the rational contraction and isolation rules of
Theorem~\ref{thm:symbol-rules}, including omission of an entire rational
prime. The same argument applies to any coefficient test with these
rules, such as the restricted even test constructed in
Section~\ref{pw:split-section}.

Encoding quadratic residue symbols by graphs has a precedent in the
congruent-number problem: Feng--Xiong express descent conditions through
prime graphs and the ranks of their binary Laplacians
\citep[\S2, Definitions~2.1 and~2.3 and Lemma~2.2]{FengXiong2004}.
Here the graphs record coefficient-symbol rules, and their oriented
contractions produce simultaneous unit values.

\subsection{Prime labels and their realization}

Let $F$ be $\mathbf Q$ or a quadratic field.  The applications use
$F=\mathbf Q(E[2])$ when this field is quadratic, and also an auxiliary
quadratic field when $E[2]$ is rational.  Fix a finite set $S$ of rational
primes containing $2$ and the ramification support of the fixed data.
Enlarge $S$ so that prime ideals above it generate the class group of $F$,
and let $S_F$ include these places and the infinite places.  Set
\[
 V_S(F)=\{\alpha\in F^\times/F^{\times2}:
       v_{\mathfrak p}(\alpha)\equiv0\pmod2
       \text{ for }\mathfrak p\notin S_F\}.
\]
The valuation sequence expresses this as an extension of the squareclasses
of $S_F$-units by a subgroup of the $2$-torsion of the $S_F$-class group;
in particular, it is finite-dimensional over $\mathbf F_2$
\citep[Theorems~4.4 and~5.11]{MilneANT2020}.
Write $F_S^{(2)}=F(\sqrt{V_S(F)})$.

A \emph{split vertex} is a rational prime $p\notin S$ split in $F$,
with an \emph{orientation} $\mathfrak p\mid p$ when $F$ is quadratic.
Choose a generator $\pi_{\mathfrak p}$ whose divisor outside $S_F$ is
$\mathfrak p$; this is possible by the class-group condition.  Use its
conjugate for the other orientation.  Then
\begin{equation}\label{gr:norm}
 N_{F/\mathbf Q}(\pi_{\mathfrak p})=p u_{\mathfrak p},
 \qquad u_{\mathfrak p}\text{ a signed rational }S\text{-unit}.
\end{equation}
For $F=\mathbf Q$, use the generator $p$.  In the quadratic case an inert
rational prime is a \emph{simple vertex}, with rational generator $p$.
A split vertex has its fixed Frobenius label on $F_S^{(2)}$; a simple
vertex has rational fixed-support labels.  A \emph{type} may also specify
local squareclasses, support valuation parities, real signs, and finite ray
conditions on its generator.  Prescribed types must be compatible and
individually realizable.  A \emph{filter} fixes the total rational labels
and sign of the signed-prime product, where $p^*=(-1/p)p$.

Write quadratic residue symbols additively in $\mathbf F_2$.  Between two
split vertices in the quadratic case there are two independent forward
bits: the residues of $\pi_{\mathfrak p}$ and
$\overline{\pi_{\mathfrak p}}$ at the chosen place of the other vertex.
Hilbert reciprocity and the types determine the reverse bits.
Their sum, with the fixed correction in \eqref{gr:norm}, gives the rational
norm bit.  If an endpoint is simple, or if $F=\mathbf Q$, there is one
rational forward bit.  Vertices, types, and these bits form an
\emph{arithmetic graph}; there is no self-residue datum.

For the coefficient symbols, the vertices are the prime factors of the
Fourier index. Their types fix the support and elementary Frobenius
conditions, while the edge bits record the remaining quadratic residue
tests. Varying an edge leaves those types and the total filter unchanged.
The next calculation explains the resulting contraction rule.

For the standard rational-two-torsion symbol of
Proposition~\ref{tc:rational-rules}, take split primes $p,q$ and let $U$
contain $p$ but not $q$, with matching
parity. Write $V=U/p$ and $g=\Fr_q$. By
Proposition~\ref{tc:rational-rules},
\[
 A(Uq)=\bigl[t(g)H_U\bigr][1].
\]
Changing one residue bit from $p$ to $q$ replaces $g$ by $g\tau_p$,
where $\tau_p$ is the corresponding inertia element. The level identity
\eqref{tc:exact-level-identity} moves the difference to the coefficient
at $p$. Its two first-shell values are
\[
 c_p(D')A(D'),\qquad (c_p(D')+1)A(D'),
 \qquad D'=Vq',\quad \Fr_{q'}\sim g\phi_p.
\]
The omitted source is unramified at the whole prime $p$, so its value
$A(D')$ is the same in the two terms. Their sum is $A(D')$.
Thus, if $z$ is the changed edge bit,
\begin{equation}\label{gr:first-contraction}
 A\big|_{z=0}+A\big|_{z=1}=A(Vq').
\end{equation}
Here $q'$ is a fresh prime carrying the product labels of $p$ and $q$;
it replaces two vertices by one. Its integer value is not the product
$pq$.

In the quadratic field, the two choices of oriented bit correspond to
using $\pi_{\mathfrak p}\pi_{\mathfrak q}$ or
$\pi_{\mathfrak p}\overline{\pi_{\mathfrak q}}$ as a product template,
up to the choice of output orientation. The template adds the support
labels and the external residue bits in the indicated orientations.
There is no internal residue attached to the new vertex. Omitting $p$
removes both places above that rational prime from the source, so a
second difference in the other oriented bit gives zero. Split--simple
contraction has the same interpretation with rational product data;
a simple--simple difference is zero.

To use these formulas on arbitrary arrays of bits, we must know that
the arrays and the product templates have prime representatives, and
that different representatives give the same symbol. The following
lemma supplies both facts.

\begin{lemma}[Finite realization and product labels]\label{gr:realization}
Fix finitely many old vertices and individually realizable new types.
Every choice of independent forward bits is realized by distinct fresh
primes, avoiding any given finite set.  Further fixed local tests may be
included in $S$.

A split cluster may be represented by the product of its oriented
generators, with conjugations prescribed along a contraction tree; its
support types and external bits are the corresponding sums.  A cluster
containing one simple vertex has the corresponding rational product data.
These products may be replaced by fresh prime representatives.  For a
symbol satisfying the omitted-vertex dependence of
Theorem~\ref{thm:symbol-rules}, the graph data determine its value, and
these replacements preserve the contracted symbol.
\end{lemma}

\begin{proof}
The displayed generators account for every relevant radical.  Let $U$
be a finite set of new rational primes and $U_F$ the places above them.
From a squareclass with even valuations outside $S_F\cup U_F$ subtract
the generators at the primes of $U_F$ where its valuation is odd.  The
remainder lies in $V_S(F)$.  At an inert prime its rational generator has
valuation one at the unique place of $F$.  Thus the fixed labels and
residue bits determine every elementary Frobenius test with this
ramification support.

For a new split vertex, start with a template $\pi_0$ of its type, at a
prime ideal $\mathfrak p_0$ away from the old vertices.  Weak approximation
gives $a\in F^\times$ sufficiently close to $1$ at $S_F$, positive at
the specified real places, and with residues at the old prime ideals
chosen so that $a\pi_0$ has the prescribed quadratic residues.  Use a
ray modulus containing deep enough powers of $S_F$, all old prime ideals,
and the real places.  A prime ideal $\mathfrak q$ in the ray class of
$(a)\mathfrak p_0$ satisfies
\[
 \mathfrak q=(ab)\mathfrak p_0,\qquad
 b\equiv1\pmod{\text{ray modulus}}.
\]
Then $ab\pi_0$ has the desired divisor and type.  The prime ideals in
this ray class have positive Dirichlet density
\citep[Chapter~V, Theorem~2.5]{MilneCFT2020}.
For $s>1$, the sum of $(N\mathfrak q)^{-s}$ over prime ideals of
residue degree greater than one is at most $2\sum_p p^{-2s}$, which
stays bounded as $s\downarrow1$.  The corresponding sum over the ray
class diverges.  Hence that class contains infinitely many degree-one
primes.  Outside the fixed ramification support these give split rational
primes, and finite exclusions remove finitely many choices.

At an old inert prime, prescribe instead the norm of the residue in
$\mathbf F_{p^2}^\times$.  Its norm onto $\mathbf F_p^\times$ is
surjective.  For a new simple vertex all conditions are rational
congruences, including the inert class in $F$; the Chinese remainder
theorem and Dirichlet's theorem apply.  Successive selection realizes the
whole graph.  Reciprocity determines reverse bits, since the support
Hilbert symbols are fixed and only the two endpoint residues remain.
Taking norms gives \eqref{gr:norm}, with no extra independent condition.

The same construction applies to a product template.  It preserves its
support data and its residues at every remaining vertex; no internal
residue is imposed on the replacement.  For a simple cluster use rational
congruences for the norm labels.  To check the rational filter in a split
cluster, write $\pi=\prod_i\pi_i$ and $\pi'=ab\pi$ for its replacement.
The multiplier $ab$ is a square unit at $S_F$ and positive at real places.
Consequently
\[
 N\pi=(\prod_i p_i)u,\qquad N\pi'=q u
\]
have the same signed rational $S$-unit $u$: its valuations and sign agree.
Thus $q/\prod_i p_i=N(ab)$ is locally square at the rational support.
The condition at $2$ also identifies $q^*$ with $\prod_i p_i^*$ in the
fixed tests.  This verifies preservation of the norm correction.

Finally, compare two disjoint realizations of the same graph.  Insert
fresh bridge vertices with matching bits to both realizations and replace
vertices one at a time.  The omitted-vertex Frobenius tests agree by the
radical description above, so the symbol values agree.  A preliminary
bridge removes any overlap.  This proves both graph dependence and
invariance under the product replacements.
\end{proof}

Consequently, fixing the vertex types gives a well-defined function
$A_{\mathcal G}$ of the independent forward edge bits: choose any prime
realization of the graph $\mathcal G$ and evaluate its coefficient
symbol. Lemma~\ref{gr:realization} makes the result independent of that
choice. We will write its unique multilinear polynomial as $P$.
Equation~\eqref{gr:first-contraction} says that a derivative of $P$
evaluates the contracted graph.

A split vertex is \emph{primary} if its generator is a square unit at
every finite place of $S_F$ and positive at every real place.  The principal
ray class realizes such types.  Reciprocity makes their fixed radical
labels trivial, and their signed rational contributions are squares in the
fixed tests.  Including stationary vertex primes in those tests makes
primary incidences to them neutral.  Including the support of a quadratic
discriminant $k$ ensures that every primary prime splits in
$\mathbf Q(\sqrt{k})$.

Two symbols may use different fixed-support tests or read only rational
norm labels from this same graph. We call these \emph{compatible
projections} when they forget fixed-label coordinates or take norm data,
preserve primary vertices, and commute with the
omitted-vertex product substitution. In the norm projection either
oriented-bit toggle changes the rational bit once. Thus one physical
edge difference induces the specified difference for each symbol; it
need not give the same polynomial for both symbols.

\begin{definition}[Witnesses and permitted partitions]\label{gr:partitionable}
A \emph{witness} is a finite graph with symbol $1$.  A graph is
\emph{partitionable} if it admits a partition into nonempty clusters, each
containing at most one simple vertex, and oriented contraction trees whose
contracted graph is a witness.  The contraction trees are part of the data;
this definition makes no assertion about a Selmer group.
\end{definition}

For a two-isogenous curve with quadratic field $F'=\mathbf Q(E'[2])$,
a simple vertex is \emph{good} if it is inert in both $F$ and $F'$.
A \emph{donor} is split in $F$ and inert in $F'$.  Attaching a donor
turns a simple vertex split in $F'$ into a good one.  These are conditions
on finite labels; their supply will be established in the arithmetic
application.

\subsection{Forests as finite-difference operators}

For an edge bit $z$, put $\Delta_z f=f|_{z=0}+f|_{z=1}$.
Every function on a finite binary cube has a unique multilinear polynomial
over $\mathbf F_2$; $\Delta_z$ is its formal derivative.  Differences
therefore commute and satisfy $\Delta_z^2=0$.

\begin{lemma}[Oriented forests]\label{gr:forest}
Assume the contraction, isolation, and omitted-prime assertions of
Theorem~\ref{thm:symbol-rules}.  Differentiation along a tree contracts
it, adding product labels and external incidences; the result is zero if
the tree contains two simple vertices.  A collection of edge variables
containing a cycle has zero derivative, including when its orientations
are incompatible.  The two oriented variables on one pair form a cycle
of length two.
\end{lemma}

\begin{proof}
The one-edge rule is the symbol contraction identity, with product
replacement supplied by Lemma~\ref{gr:realization}.  Induction contracts
a tree: split--split mergers remain split, split--simple mergers remain
simple, and merging two simple components gives zero.  A root orientation
determines all relative orientations along the tree.  Its product template
uses support types and external bits, not the differentiated internal bits.

For the two variables on a split pair $p,q$, differentiate first in one
inertia bit over $p$.  The resulting symbol omits the \emph{entire rational
prime} $p$, so its elementary extension is unramified at both places over
$p$.  The derivative in the other oriented bit is zero.  Repeated use of
the same variable gives zero by $\Delta_z^2=0$.

For a longer cycle, contract a spanning path.  The closing edge becomes
internal to one cluster and changes neither its support label nor an
incidence to a remaining vertex.  The omitted-endpoint argument just given
therefore kills its derivative.  Conjugating the product only changes its
orientation; it does not restore the omitted endpoint or create a
self-residue.

For example, after contracting aligned edge $12$, the two bits from its
product to vertex $3$ are
$t_0=z_{13}^0+z_{23}^0$ and $t_1=z_{13}^1+z_{23}^1$, up to fixed type
corrections.  An aligned edge $23$ differentiates $t_0$.  The closing
edge $31$ differentiates either $t_0$ again, giving zero, or $t_1$, giving
the omitted-endpoint zero.  Every longer cycle reduces to these cases.
Further differences commute with this zero derivative.
\end{proof}

Call the original vertices \emph{terminals}; any stationary external
vertices retain their types and incidences. Fix their vertex types and a
realizable baseline graph. Write $z=(z_e)$ for the independent toggles of
the terminal--terminal bits, and let $P(z)$ be the terminal symbol.
A \emph{network} adjoins primary vertices and specified incidences
involving them. It adds no terminal--terminal edge: the original bits
$z$ remain parameters. All unspecified primary incidences are neutral.
Within each network fix a consistent orientation at each terminal.
Different networks may use different alignments.

Write $y=(y_e)$ for the bits on the added incidences, and let $f(z,y)$
be the symbol of the enlarged graph. We shall express $f(z,\mathbf1)$
as a finite-difference operator applied to $P(z)$. The differences in
$y$ first contract the added graph; the forest rule then expresses each
surviving term as differences in the original variables $z$.

For example, join one primary vertex $r$ to two permitted terminals
$a,b$, in a fixed alignment. Expanding in its two added incidences gives
four terms. The term with no difference is zero because $r$ is isolated.
Each one-edge difference contracts $r$ onto one terminal and leaves the
original symbol $P$, so these two terms cancel in $\F_2$. The two-edge
tree contracts the terminal pair and gives $\Delta_{z_{ab}}P$.
Thus this path acts exactly as the terminal difference
$\Delta_{z_{ab}}$. The path has two added edges, but its operator has
contraction degree one. In general, a component meeting $t$ terminals
has terminal contraction degree $t-1$.

The lowest two operator degrees can be read from a matrix. Index the
internal primary vertices by $I$, let $e_{ij}$ be their internal edge
bits, and let $c_a$ be the incidence column from them to terminal $a$.
The grounded Laplacian is the symmetric matrix over $\F_2$ with
\[
 Q_{ij}=e_{ij}\quad(i\ne j),\qquad
 Q_{ii}=\sum_{j\ne i}e_{ij}+\sum_a(c_a)_i.
\]
Equivalently,
\begin{equation}\label{gr:grounding-dictionary}
                         Q\mathbf1=\sum_a c_a.
\end{equation}
Thus prescribed off-diagonal entries and terminal columns realize a
given symmetric $Q$ precisely when this row-sum identity holds. For the
one-primary-vertex path above, $Q=(0)$ and $c_a=c_b=(1)$.

The coefficients below are instances over $\F_2$ of the all-minors
matrix-tree theorem \citep[\S1, Equation~(2)]{Chaiken1982}; we include
the incidence-matrix proof needed here.

\begin{lemma}[Network expansion]\label{gr:network}
A network satisfies $f(z,\mathbf1)=\mathcal N P(z)$, where
$\mathcal N$ is a sum of terminal finite-difference operators indexed by
forests. In each surviving forest term, every component containing an
internal primary vertex meets a terminal. A component meeting $t$
terminals has contraction degree $t-1$;
a component containing two simple terminals acts as zero.  Networks with
disjoint internal vertices and no edges between them multiply as operators
on the same $P$.  In particular, degree greater than $\deg P$ annihilates it.

For the grounded Laplacian $Q$ and terminal columns $c_a$ defined above,
the constant
and the pair-contraction coefficients are
\begin{equation}\label{gr:forest-coefficients}
 \det Q,\qquad c_a^{\mathsf t}\operatorname{adj}(Q)c_b\quad(a\ne b).
\end{equation}
The latter multiplies the terminal difference in the selected alignment;
a prohibited pair multiplies the zero operator.
\end{lemma}

\begin{proof}
For the set $E_0$ of added incidences, keep the terminal variables $z$
fixed and expand in the variables $y$:
\[
 f(z,\mathbf1)=\sum_{T\subseteq E_0}(\Delta_T^y f)(z,\mathbf0).
\]
Lemma~\ref{gr:forest} removes cycles.  A component meeting no terminal
contracts to an isolated primary vertex, so isolation removes it, including
an untouched primary vertex.  In a surviving component the primary labels
and external primary bits are trivial.  Its contracted label and external
incidences are thus exactly those obtained by contracting its terminals
in the original graph.  A terminal tree with the same relative orientations
gives that operator, of degree $t-1$.

Expanding two disjoint networks gives composition: product labels agree
in either order, and any resulting cycle vanishes.  Polynomial
differentiation then proves the degree bound.  This reasoning also verifies
compatibility with projections at every degree.  Norms commute with
oriented products, primary norm corrections are neutral, and either
oriented toggle induces its rational toggle on the original terminal
polynomial.  Repeated projected differences or prohibited simple mergers
give zero.  Forgetting fixed-label coordinates commutes with the same
products.

For the coefficients, write the Laplacian of the added network as
$L=DD^{\mathsf t}$,
where each incidence column has ones at its edge endpoints.  Let $I$ and
$A$ be the internal and terminal sets.  Cauchy--Binet for $Q=L_{I,I}$
counts edge sets whose incidence matrix becomes invertible after deleting
the terminal rows.  They are exactly forests with one terminal in each
component: cycles give dependent columns, a component with no deleted row
gives dependent rows, and a tree with one row deleted is invertible by
leaf removal.  These are the degree-zero terms.

Apply Cauchy--Binet instead to row set $I\cup\{a\}$ and column set
$I\cup\{b\}$.  Each component must contain one root from each of
$A\setminus\{a\}$ and $A\setminus\{b\}$.  Thus $a,b$ lie together
and every other terminal lies in its own component.  The corresponding
minor is
\[
 \det\begin{pmatrix}Q&c_b\\c_a^{\mathsf t}&0\end{pmatrix}
   =c_a^{\mathsf t}\operatorname{adj}(Q)c_b.
\]
The corner is zero because there is no added terminal--terminal edge.  Signs
disappear in characteristic two, and the adjugate identity remains valid
when $Q$ is singular.
\end{proof}

\subsection{Selecting one address}

An \emph{address} is $x\in\mathbf F_2^b$, where $b\ge1$.
A vertex with activation $\lambda\in(\mathbf F_2^b)^*$ is present
when $\lambda(x)=1$.  Fix a nonzero target $x_*$ and an edge bit $z$
joining terminals $a_0,a_1$ active there.  Either all terminals are always
present, or all have nonzero linear activations with at least one terminal
for each nonzero linear form.  For one symbol we need a derivative as the leading term at $x_*$ and
a constant term $1$ elsewhere. For two symbols we must also account for
a derivative allocated to just one factor; the second network includes
a binary parameter for this purpose. We construct these two networks by
prescribing their constant and degree-one terms; the final degree count
will remove every higher term.

Choose a dual basis $\lambda_1,\ldots,\lambda_b$ with
$\lambda_i(x_*)=1$ for every $i$.  Indeed, start with
$\lambda_1(x_*)=1$ and add it to a basis of the annihilator of $x_*$.
Make $b$ groups of three primary vertices, with all three vertices in
group $i$ activated by $\lambda_i$. The physical coordinate vectors index
these actual vertices. In each three-dimensional coordinate space choose
a computational basis $(e_i,u_i,v_i)$ with $e_i=(1,1,1)$. Prescribe a
symmetric form $Q$ by making the $e_i$ orthonormal and orthogonal to all
$u_j,v_j$, and using on the remaining coordinates
\begin{equation}\label{gr:H}
 H(T,B)=\begin{pmatrix}0&T\\T^{\mathsf t}&B\end{pmatrix}.
\end{equation}
Convert this form to its matrix in the physical vertex coordinates,
which we also denote by $Q$. Write $Q_J$ for the restriction to the
physical coordinates in groups $J$.

The column $h_i=Qe_i$ represents the covector dual to $e_i$. Since the
coordinate changes do not mix groups, $h_i$ is supported on group $i$;
it is independent of $T,B$. Attach this incidence column to a terminal
of activation $\lambda_i$, or to any fixed terminal in the stationary
case. When the active groups are $J$, the sum of these anchor columns is
\[
             \sum_{i\in J}h_i|_J=Q_J\mathbf1_J,
\]
since $\mathbf1_J=\sum_{i\in J}e_i$. Use the off-diagonal physical
entries of $Q$ as the internal edge bits. Equation~\eqref{gr:grounding-dictionary}
then gives the prescribed diagonal as well: the anchors make the actual
grounded Laplacian equal to $Q_J$ at every address, even after groups
have disappeared.

Also attach the same column $c$ to $a_0,a_1$.  If their moving activations
$\alpha,\beta$ differ, attach it to a terminal of activation
$\alpha+\beta$ as well.  The additional column sum is always zero:
\[
 \bigl(\alpha(x)+\beta(x)+(\alpha+\beta)(x)\bigr)c=0.
\]
The third terminal is absent at $x_*$.  Equal activations, or two stationary
terminals, require only the first two columns.  Coincident terminal roles
mean addition of columns.  Hence these attachments preserve the grounding
at every restriction.

Write $O_{\ge d}$ for a sum of contraction operators of degrees at least
$d$.  This denotes a filtration, not a numerical estimate.
Figure~\ref{gr:selector-figure} depicts the group-level cycle and path
used in the two programs when $b=4$.

\begin{figure}[tbp]
\centering
\setlength{\unitlength}{1mm}
\begin{picture}(115,36)
\put(15,28){\circle{6}}\put(15,28){\makebox(0,0){$1$}}
\put(40,28){\circle{6}}\put(40,28){\makebox(0,0){$2$}}
\put(40,8){\circle{6}}\put(40,8){\makebox(0,0){$3$}}
\put(15,8){\circle{6}}\put(15,8){\makebox(0,0){$4$}}
\put(18,28){\vector(1,0){19}}\put(40,25){\vector(0,-1){14}}
\put(37,8){\vector(-1,0){19}}\put(15,11){\vector(0,1){14}}
\put(27.5,0){\makebox(0,0){Cycle}}
\put(75,28){\circle{6}}\put(75,28){\makebox(0,0){$1$}}
\put(100,28){\circle{6}}\put(100,28){\makebox(0,0){$2$}}
\put(100,8){\circle{6}}\put(100,8){\makebox(0,0){$3$}}
\put(75,8){\circle{6}}\put(75,8){\makebox(0,0){$4$}}
\put(78,28){\vector(1,0){19}}\put(100,25){\vector(0,-1){14}}
\put(97,8){\vector(-1,0){19}}
\put(87.5,0){\makebox(0,0){Path}}
\end{picture}
\caption{The off-diagonal patterns of $T-I$ for the two programs when
$b=4$.  The vertices represent groups, not the individual primary
vertices of the network.  A proper group restriction breaks the cycle;
the path from $1$ to $4$ survives only when every group is present.}
\label{gr:selector-figure}
\end{figure}

\begin{lemma}[Two address programs]\label{gr:programs}
With these activations and grounding, there are the following networks:
\begin{enumerate}
\item A singular program $\mathcal S$ such that
\[
 \mathcal S_{x_*}=\Delta_z+O_{\ge2},\qquad
 \mathcal S_x=1+O_{\ge1}\quad(x\ne x_*).
\]
\item Two nonsingular programs on the same sites, indexed by
$\varepsilon\in\mathbf F_2$, such that
\[
 \mathcal N_x^{(\varepsilon)}=1+O_{\ge1},\qquad
 \mathcal N_x^{(1)}+\mathcal N_x^{(0)}
 =\begin{cases}\Delta_z+O_{\ge2},&x=x_*,\\
 O_{\ge2},&x\ne x_*.
 \end{cases}
\]
\end{enumerate}
All new sites have nonzero activation and are absent at $0$.
Unspecified incidences, including those to stationary external vertices,
may remain neutral.
\end{lemma}

\begin{proof}
For the first program, take $T=I+P$, with $P$ the cyclic permutation
matrix, and let $B$ be diagonal with exactly one nonzero entry.  Every
proper principal restriction of $P$ is a union of directed paths, hence
nilpotent.  Thus $T_J$ and $H(T_J,B_J)$ are invertible for proper $J$,
including the empty restriction, whose determinant is $1$.

On all groups, $\ker T$ and $\ker T^{\mathsf t}$ are spanned by
$\mathbf1$.  A radical vector of $H$ has $v=t\mathbf1$; summing
$T^{\mathsf t}u+Bv=0$ gives $t=0$, since
$\mathbf1^{\mathsf t}B\mathbf1=1$.  Hence $Q$ has radical generated
by $r=\sum_i u_i$.  Its adjugate is $rr^{\mathsf t}$: it is nonzero
of rank one, symmetry identifies its factors with the kernel, and the
only nonzero scalar in $\mathbf F_2$ is $1$.

Take $c$ dual to $u_1$, so $c^{\mathsf t}r=1$.  All anchors annihilate
$r$.  At $x_*$ the only terminal columns pairing nontrivially with $r$
are therefore those of $a_0,a_1$.  Formula~\eqref{gr:forest-coefficients}
gives constant zero and precisely $\Delta_z$ in degree one.  Every other
address has a proper group subset and constant $1$, since the activation
forms form a basis.  For $b=1$, $T=0$ and $B=(1)$ give the same radical.

For the second program take $T=I+L$, where $L$ is the consecutive lower
shift, and $B_\varepsilon=\varepsilon E_{bb}$.  Every group restriction
of $T$ is invertible, and
\begin{equation}\label{gr:inverse}
 H(T,B)^{-1}=
 \begin{pmatrix}T^{-\mathsf t}BT^{-1}&T^{-\mathsf t}\\
 T^{-1}&0\end{pmatrix}.
\end{equation}
Again take $c$ dual to $u_1$.  On groups $J$, the change in
$c^{\mathsf t}Q_J^{-1}c$ is $((T_J^{-1})_{b1})^2$, interpreted as
zero if an endpoint is absent.  The expansion
$(I+L_J)^{-1}=I+L_J+L_J^2+\cdots$ shows that this is $1$ exactly
when the entire path $1,\ldots,b$ survives, namely at $x_*$.
Pairings with anchors do not change, since $Q_J^{-1}h_i=e_i$ and
$e_i,h_i,c$ are independent of $\varepsilon$.  At the target the
balancing terminal is absent, leaving only the changed pair $a_0,a_1$.
All determinants are $1$.  Lemma~\ref{gr:network} proves the assertion,
also for the single-vertex path $b=1$.
\end{proof}

\subsection{Simultaneous units on a punctured cube}

The preceding programs control only two degrees of their operators.
To make this sufficient, at each address we spend exactly the degree
of a chosen top monomial.  Any additional contraction then vanishes.
For two symbols, the same accounting uses the ordinary product of their
top homogeneous parts. The multiplication is performed before imposing
the Boolean relation $z^2=z$, so one occurrence of $z$ in each factor
remains two contraction requirements.

\begin{theorem}[Address lemma]\label{gr:address}
Assume the contraction, isolation, and whole-prime omission rules of
Proposition~\ref{tc:rational-rules}, or the same rules for another
symbol, together with the realization freedoms of
Lemma~\ref{gr:realization}.  Let $b\ge1$ and
consider either of the following terminal arrangements:
\begin{enumerate}
\item One symbol on a nonempty set of always active terminals, whose
polynomial in the independent pair bits is nonzero.
\item One or two symbols on terminals with nonzero linear activations,
including a terminal for every nonzero linear form.  At each $x\ne0$,
each symbol separately is a nonzero polynomial in the active bits.
The symbols may have stationary external data and may read compatible
projections of a common graph.
\end{enumerate}
At every $x\ne0$, assume that the active terminal data and the stationary
external data obey each required filter.  Then finitely many primary
auxiliaries with nonzero linear activations, and one realizable assignment
of graph bits, make all the symbols $1$ at every $x\ne0$.  No auxiliary
is present at $0$.  All may satisfy fixed splitting and neutral-incidence
conditions included in their primary type.

A nonzero evaluation or a nonzero iterated edge difference suffices for
the polynomial hypothesis.  In particular, a partition as in
Definition~\ref{gr:partitionable} suffices in the stationary case.
\end{theorem}

\begin{proof}
Fix a common realizable baseline for the terminal bits and regard their
toggles as variables; omit variables invisible to a symbol.  A nonzero
iterated difference implies a nonzero polynomial, while every nonzero
polynomial has a top monomial with coefficient $1$.  By
Lemma~\ref{gr:forest}, this also proves the assertion about partitions.

For one symbol, let $P_x$ have degree $d_x$ at $x\ne0$ and choose a
monomial of degree $d_x$ with coefficient $1$.  For each of its variables
adjoin a separate singular program targeted at $x$, all with disjoint
internal vertices.  At address $x$, its own $d_x$ programs have zero
constant, so they consume at least $d_x$ contraction degrees.  Any extra
degree from any program kills $P_x$.  Only the selected degree-one terms
of these programs and the constants of the others survive, giving the
chosen coefficient $1$.  A nonzero constant polynomial needs no program.
Stationary anchors give the same argument for always active terminals.

For two symbols, write their polynomials at $x$ as $P_x,R_x$, of degrees
$d_x,e_x$.  Multiply their top homogeneous parts in the ordinary ring
$\mathbf F_2[z_1,\ldots,z_m]$, before imposing Boolean relations.
This ring is a domain, so choose a coefficient $1$ at a monomial
\begin{equation}\label{gr:chosen-monomial}
 M_x=\prod_z z^{\nu_{xz}},\qquad
 \nu_{xz}\in\{0,1,2\},\qquad
 \sum_z\nu_{xz}=d_x+e_x.
\end{equation}
For each exponent two add one singular program targeted at $x$; for
each exponent one add one nonsingular program, with its own parameter
$\varepsilon_{xz}\in\mathbf F_2$.  Keep all internal sets disjoint.
Let
\[
 \Phi(\varepsilon)=\prod_{x\ne0}
    (\text{final value of }P_x)(\text{final value of }R_x).
\]
We show that $\sum_\varepsilon\Phi(\varepsilon)=1$, which implies
that one choice makes every factor $1$.

Expand the operators by contraction degree.  At its target a singular
program has zero constant in both factors, so it consumes at least one
degree in each.  A nonsingular program has constant $1$ independent of
its parameter.  Summing that parameter cancels a term using only its
constants everywhere; a surviving term must consume at least one degree
somewhere for that program.  The total of these minimum costs is
\begin{equation}\label{gr:budget}
 D=\sum_{x\ne0}(d_x+e_x).
\end{equation}
A nonzero term cannot consume more, since some factor would then be
differentiated beyond its degree.  Every surviving term therefore uses
each minimum cost exactly.

Thus a singular program contributes precisely $\Delta_z$ in each factor
at its target and constants elsewhere.  A nonsingular program contributes
degree one in exactly one factor.  Summing its parameter kills that
contribution away from its target and gives $\Delta_z$ at the target;
every other occurrence contributes the common constant $1$.  In
particular, no higher forest coefficient remains.

At one address put $Z_i=\{z:\nu_{xz}=i\}$ for $i=1,2$, and write
$z_J=\prod_{z\in J}z$.  The surviving allocation sum is
\[
 \sum_{J\subseteq Z_1}
 [z_{Z_2}z_J]P_x^{\rm top}\,
 [z_{Z_2}z_{Z_1\setminus J}]R_x^{\rm top}
   =[M_x](P_x^{\rm top}R_x^{\rm top})=1.
\]
Wrong-degree coefficients are zero.  Since the total degree is $D$
and no factor exceeds its degree, each factor receives its full degree;
its derivative is exactly the displayed coefficient, independent of the
baseline.  Allocations at distinct addresses are independent, so their
product is $1$.  This covers constant factors too.  For instance, if
both top forms are $z$, their ordinary product is $z^2$: a singular
program differentiates each once.  Boolean reduction would lose that
allocation.

Choose parameters with $\Phi=1$ and realize the resulting finite graph
by Lemma~\ref{gr:realization}.  Primary types preserve the filters and
the stipulated neutral incidences and splitting conditions.  Nonzero
linear activations ensure that no new vertex appears at $0$.
\end{proof}

For a stabilized symbol this is a finite construction at each dimension.
First choose $b$, the terminal data, and all networks; their weights and
coefficient tests then have a finite bound.  Select an actual stage where the numerical bounds and the finitely many stabilized symbol
evaluations hold for that range.  The required precision may
grow with $b$.  The unit-valuation constant, however, is the constant of
the coefficient construction and must be independent of $b$; the forest
argument supplies no additional uniformity assertion.
